\addvspace{1cm}
\subsection{Rechenregeln der z-Transformation}
\label{sec:Sätze_der_ZT}
Genauso wie die Laplace-Transformation, gelten auch für die z-Transformation einige Regeln, die das Rechnen mit der z-Transformierten erleichtern. Einige dieser Regeln sind ähnlich, während manche Sätze sich stark durch die zeitliche Diskretisierung unterscheiden.
\addvspace{1cm}
\textbf{Linearität}
Da es sich bei der z-Transformation um eine lineare Transformation handelt, gilt für zwei kausale, zeitdiskrete Abtastfolgen $x_1[k]$ und $x_2[k]$ mit den zugehörigen Bildfunktionen $X_1(z)$ und $X_2(z)$ der Zusammenhang:
\begin{equation}
    a \cdot x_1[k] + b \cdot x_2[k] \TransformHoriz a \cdot X_1(z) + b \cdot X_2(z)
\end{equation}

Die Linearität ermöglicht es, zusammengesetzte Funktionen einfach zu transformieren, da die Bestandteile einzeln behandelt werden können.

\addvspace{1cm}
\textbf{Verschiebungssatz}
\label{sec:Verschiebungssatz}

Wird die Abtastfolge $x[k]$ um $k_0$ Schritte nach rechts verschoben, so gilt auf Grund der der Definition der z-Transformation::
\begin{equation}
    x[k-k_0] \TransformHoriz z^{-k_0} \cdot X(z)
\end{equation}

\addvspace{1cm}
\textbf{Dämpfungssatz}

Wird die Abtastfolge $x[k]$ mit einem Faktor $a^k$ gedämpft, so gilt:

\begin{equation}
    x[k] \cdot a^k \TransformHoriz X(\frac{z}{a})
\end{equation}

Bei der Dämpfung ist zu beachten, dass sich der Konvergenzbereich ändert. Dieser hängt von dem Dämpfungsfaktor $a$ ab. Der o.g. Zusammenhang gilt nur für $|\frac{z}{a}| > 1$.
\addvspace{1cm}
\textbf{Multiplikation im Zeitbereich}

Wird jeder Abtastwert der Folge $x[k]$ mit dem Wert von $k$ gewichtet, so entspricht dies der Ableitung der Bildfunktion mit dem Vorfaktor $-z$:
\begin{equation}
    x[k] \cdot k \TransformHoriz -z \cdot \frac{d}{dz} X(z)
\end{equation}

Über diesen Zusammenhang lässt sich sehr leicht zeigen, wie die Korrespondenz der mit k gewichteten Sprungfolge aus Tabelle \ref{tab:ZT_Korrespondenzen} zustande kommt:

\begin{equation}
    \mathcal{Z}(k \cdot x[k]) = -z \cdot \frac{d}{dz} \left(\frac{z}{z-1}\right) = -z \cdot \left(\frac{(z-1) - z}{(z-1)^2}\right) = -z \cdot \left(\frac{-1}{(z-1)^2}\right) = \frac{z}{(z-1)^2}
\end{equation}
\addvspace{1cm}
\textbf{Faltungssatz}

Für die diskrete Faltung zweier Abtastfolgen gilt der Faltungssatz der z-Transformation. Er sagt aus, dass die Faltung der zwei Abtastfolgen $x_1[k]$ und $x_2[k]$ im Zeitbereich der Multiplikation der zugehörigen Bildfunktionen im Frequenzbereich entspricht:
\begin{equation}
    x_1[k] * x_2[k] \TransformHoriz X_1(z) \cdot X_2(z)
\end{equation}

Diese Regel vereinfacht stark das Rechnen mit Abtastfolgen und ist ein wichtiger Grund, die z-Transformation zu beherrschen.


\addvspace{1cm}
\textbf{Anfangswertsatz}

Für eine rechtsseitige Folge, wie sie bei Abtastwerten vorliegt, also bei der $x[k] = 0$ für $k<0$ gilt, ist laut dem Anfangswertsatz der Wert zum Zeitpunkt $k=0$ :
\begin{equation}
    x[0] = \lim_{z \to \infty} X(z)
\end{equation}

\cite[S.215]{book3} \cite[S.239]{book2}